\chapter{Methodology and Evidence Standards}
\label{chap:methodology}

\classification{Evidence status: Methodological chapter. Confidence baseline: High. This chapter documents the investigation process, repository controls, and publication standards. It does not assess case hypotheses or introduce new evidence.}

\section{Objectives}

The objective of this investigation is to produce a reproducible forensic and intelligence-style dossier on the unknown man found at Somerton Beach. The project separates evidence acquisition, claim verification, chronology reconstruction, forensic assessment, physical-evidence cataloguing, contextual analysis, and manuscript production into distinct controlled stages.

The scope is limited to documented evidence, structured analysis, and clearly identified intelligence gaps. Unsupported prior claims, speculative scenarios, and unsourced narrative elements are retained only as research history or hypothesis material. They are not promoted into factual findings unless supported by the repository evidence base.

The research philosophy is evidence-first and audit-led. The repository is treated as the primary asset; the manuscript is an output of that database. A future researcher should be able to trace a published statement back to a claim, source, document, event, exhibit, person, location, register, or decision-log entry.

The intended contribution is therefore methodological as well as historical. The dossier is designed to preserve uncertainty, expose source limitations, and make later correction possible if new records or forensic material become available.

\section{Intelligence Dossier Methodology}

An intelligence-assessment model was selected because the case contains incomplete records, contested interpretations, physical-evidence custody gaps, later media claims, and historically plausible but weakly evidenced contextual theories. A traditional narrative history can obscure those distinctions by arranging uncertain material into a continuous story. This dossier instead separates what is documented, what is inferred, what is hypothesized, and what remains unknown.

The method is evidence-first. Sources are acquired or catalogued before claims are promoted. Claims are then checked against the strongest available evidence, with primary and official records preferred where available. Where a primary record exists, later summaries do not replace it.

Reproducibility is supported through machine-readable registers, stable identifiers, source metadata, preservation manifests, page-level indexes, evidence maps, and chapter validation reports. The chapter text is not the evidence system; it is a controlled synthesis of the evidence system.

Transparency is maintained through explicit source classes, confidence language, decision logging, and publication-readiness audits. Methodological decisions that affect future wording are recorded in the decision log rather than left as undocumented editorial preference.

Uncertainty preservation is mandatory. Conflicts are recorded, not smoothed. Negative findings are limited to the corpus actually reviewed. Contextual plausibility is not treated as case-specific proof.

\section{Evidence Hierarchy}

The source hierarchy determines how evidence is weighted and how conflicts are handled. It is a weighting system, not a guarantee that every higher-ranked source is complete or error-free.

\begin{table}[htbp]
\centering
\caption{Evidence hierarchy used for manuscript drafting}
\label{tab:methodology-evidence-hierarchy}
\begin{tabularx}{\textwidth}{p{0.24\textwidth}X p{0.22\textwidth}}
\toprule
Source class & Use in the dossier & Primary limitation \\
\midrule
Original official records & Highest-weight evidence for factual reconstruction, legal process, medical observations, property handling, and official actions. & Legibility, completeness, access restrictions, and missing enclosures. \\
Certified copies and official scans & High-weight substitutes where original custody is documented or an official reproduction is available. & Copy quality and possible omission of attached materials. \\
Contemporary official statements and legal records & Used for current institutional position, public legal notices, and official-process status. & Limited to what the issuing body stated publicly. \\
Contemporary newspapers & Used for chronology leads, public reporting, and contemporaneous public record where primary records are missing. & Press error, OCR distortion, compression of testimony, and sensational framing. \\
Research publications and specialist accounts & Used for technical context, source discovery, and specialist interpretation. & May not have access to full primary files or may inherit earlier errors. \\
Media reports and later synthesis & Used for OSINT updates, public claims, interviews, and leads requiring source separation. & Often mix fact, expert opinion, and interpretation. \\
Contextual material & Used to explain historical environment, institutions, capability, and limitations. & Does not establish direct case linkage without a case-specific source. \\
Internal prior analysis & Used only to extract claims, questions, and hypotheses for verification. & Not evidentiary authority. \\
\bottomrule
\end{tabularx}
\end{table}

Primary sources include coronial records, police files, pathology and toxicology material, official photographs, evidence inventories, property records, government records, and contemporary witness or legal material. Secondary sources include academic publications, books by recognized researchers, scientific commentary, and specialist reconstructions. Media reports and OSINT sources are used cautiously and must be separated from their underlying original sources when possible.

\section{Verification Framework}

Verification begins with claim extraction. Claims from the prior analytical material and later source leads are assigned stable identifiers and classified by topic. Each claim is then mapped to possible source material before it is assessed.

The verification matrix records the original claim, verification status, supporting documents, contradictory documents, page references where available, evidence summary, confidence, outstanding questions, and recommended revision. The aim is not to force a claim into true or false form, but to state how much of the wording the evidence actually supports.

Verification status values distinguish between verified, partially verified, contradicted, unsupported, unresolved, and unable to verify. These categories prevent unsupported claims from entering the manuscript as established facts and prevent partially supported claims from being overstated.

Confidence assignment reflects evidence quality. It does not express the probability of a theory. A claim supported by an authoritative primary record may receive high confidence. A claim supported only by later reporting, OCR-distorted text, or an incomplete reproduction receives lower confidence or a qualified status.

Contradiction handling is explicit. If sources disagree, both the supporting and limiting records are retained in the conflict register or relevant evidence matrix. The manuscript must describe the disagreement or limitation rather than resolve it by editorial assumption.

Unsupported claims remain part of the research record but are downgraded, excluded, or moved into controlled hypothesis or historiography discussion. The prior analytical dossier is therefore preserved as a source of questions and claims, not as case evidence.

\section{Research Database}

The repository is organized as a structured research database. Its registers are designed to be searchable, expandable, and auditable. Stable identifiers allow later chapters to cite evidence without renumbering or changing the underlying data model.

The evidence database contains source, evidence, person, location, timeline, intelligence-gap, hypothesis, image, research-question, and decision-log registers. These records provide the base layer for later verification and manuscript drafting.

The primary-source layer records acquired documents, archive details, reference numbers, authenticity status, preservation metadata, checksums where available, access restrictions, and document-to-claim mappings. It also preserves unavailable or restricted records so that gaps are visible rather than silently omitted.

The chronology layer records dated events with date precision, time precision, related people, related evidence, source references, confidence, and uncertainty flags. It is designed to prevent uncertain events from being merged into a smoother but less defensible timeline.

The witness and exhibit indexes provide page-level document indexing, witness/declarant records, exhibit/property records, transcription queues, and cross-reference maps. Their function is to make testimonial and exhibit material traceable before it is summarized.

The victimology layer records the known and unknown elements of identity, movement, possessions, behaviour, and relationships. It separates documented observations from biographical leads and unknowns.

The forensic-medical layer records medical observations, toxicology issues, time-of-death statements, forensic identifiers, medical uncertainties, and priority transcriptions. It separates observed pathology, testing, interpretation, and unresolved medical questions.

The physical-evidence layer records museum-style object entries, scene-reconstruction components, custody confidence, preservation history, source mapping, and unresolved physical-evidence questions. Its purpose is to distinguish object existence from object interpretation and custody strength.

The intelligence layer records capability assessments, an analysis of competing hypotheses, key assumptions, indicators and signposts, source reliability, evidence weighting, confidence assessments, and intelligence requirements. It is used to structure analysis without converting hypotheses into conclusions.

The OSINT layer records post-archive public-source updates, source reliability, conflicts with earlier assessment, new research questions, and recommended register updates. It is used as an update layer, not as a substitute for official records where those records are required.

\section{Structured Analytic Techniques}

Structured analytic techniques are used to discipline inference. They are not used to produce a preferred theory in advance of the evidence.

Analysis of Competing Hypotheses organizes possible explanatory models against supporting evidence, contradictory or limiting evidence, critical missing evidence, and current confidence. The method is designed to identify diagnostic evidence and weak assumptions.

The key assumptions check identifies assumptions that could distort later interpretation, including assumptions about identity continuity, label removal, toxicology, code interpretation, historical context, and official status. Assumptions are treated as risks to be managed rather than as hidden premises.

Capability assessment distinguishes what was historically or technically possible from what is demonstrated in this case. This prevents a capability, such as a difficult-to-detect poison or intelligence activity in the period, from being treated as evidence that the capability was used.

The evidence weight matrix records the relative strength and limitation of evidence clusters. Medical observations, property and clothing observations, witness timelines, Rubaiyat material, modern identification material, and Cold War context do not carry the same evidentiary weight.

The confidence framework records major assessment areas and the reasons for their confidence levels. Confidence is tied to source quality, completeness, custody, and diagnostic value. It is not a rhetorical device.

Intelligence requirements record the questions that would materially improve the assessment if answered. They identify missing evidence, priority, responsible chapter, and status. They also prevent unresolved questions from being quietly absorbed into narrative prose.

\section{Limitations}

The investigation is limited by unavailable and restricted records. Official custodian copies, police property records, complete forensic laboratory documentation, image rights, and some archival files remain outside the acquired record or require separate access procedures.

The archive also contains later reproductions and OCR-dependent files. OCR is useful for searching, but it is not sufficient for exact quotation where the page image is distorted or technical language is unclear. Exact wording should be checked visually before it carries interpretive weight.

Copyright and archive restrictions limit reproduction. Some documents and images can be described, catalogued, or cited but should not be reproduced unless rights, licence, or permission status is resolved. Figure insertion is therefore separated from evidence discussion.

Historical limitations are inherent to the case. Records were created under the investigative, scientific, administrative, and preservation practices of their time. Later absence of material may reflect non-retention, loss, restriction, or incomplete acquisition rather than evidentiary insignificance.

Research limitations are recorded as intelligence gaps rather than hidden. A missing source, unavailable test, uncertain custody trail, or incomplete transcription is treated as a constraint on wording.

\section{Ethics}

Objectivity is the governing ethical requirement. The manuscript must not use evidence to advocate for a preferred theory, imply guilt or motive beyond the record, or convert ambiguity into narrative force.

Source transparency is an ethical requirement as well as a methodological one. Readers should be able to distinguish primary evidence, later reporting, researcher interpretation, contextual material, and internal project analysis.

Historical responsibility requires restraint. The case involves deceased and possibly living or recently living individuals, family members, witnesses, public officials, and researchers. The manuscript should avoid sensational phrasing, unnecessary personal speculation, and unsupported psychological or moral judgments.

Reproducibility is part of research integrity. Registers, source metadata, evidence maps, decision logs, and validation reports make it possible for a later researcher to audit or challenge the work without relying on narrative trust.

\section{Repository Governance}

Repository governance is maintained through version control, stable identifiers, decision logs, publication-readiness controls, and chapter validation reports. The database is not to be silently rewritten to fit chapter prose.

Decision logging records methodological choices that affect the manuscript. This includes treatment of prior analysis, identity-status wording, publication freeze rules, bibliography integration, chapter dependency order, and authoring-framework requirements.

The research freeze established the repository as ready for manuscript drafting, with restrictions. Figure insertion remains deferred until rights are resolved. Webb identity wording remains conditional until official records are acquired. Future records may update the database, but those updates must be logged rather than absorbed silently into prose.

Future releases should preserve the evidence database, verification matrix, chronology, registers, bibliography, methodology, and release notes together. A versioned release model allows the dossier to remain corrigible while preserving the evidentiary basis for each published version.

\section{Chapter Completion Standard}

A chapter is complete only when its factual statements map to evidence, its terminology follows the controlled vocabulary, its tables and cross-references compile, its figures are rights-cleared or deferred, and its unresolved issues are stated openly. No chapter may contain draft filler text, unsupported claims, unresolved references, or unlogged methodological departures.

This standard applies to every chapter after this one. The methodology chapter therefore functions as both explanation and control: it tells the reader how the investigation was conducted and tells future drafting work how the manuscript must be built.
